% GENERATED FILE. Edit canonical lessons, then run npm run generate:books.
% canonical-source-hash: fnv1a64:77c55b2ed839b11d
% canonical-lessons: SA-C536-adrstam, SA-C536-bhavitavyata, SA-C536-samyogah, SA-C536-yadrcchaya, SA-C536-sampadyate

\chapter{A Stroke of Luck}
\label{ch:sa-a-stroke-of-luck}
\hlchaptermodality{536}

\begin{chapteropening}
\textbf{By the end of this chapter:} \emph{I can talk about luck and what happens by chance.}

Complete the last lesson of chapter 536: I can talk about luck and what happens by chance.
\end{chapteropening}

\section[\texorpdfstring{\sk{अदृष्टम्} (adṛṣṭam)}{अदृष्टम् (adṛṣṭam)}]{\sk{अदृष्टम्} (adṛṣṭam) — the unseen; fate}

\label{lesson:SA-C536-adrstam}



\begin{quote}
Before the new one: say the Sanskrit for a breaking, then the Sanskrit for one's word.
\end{quote}

\subsection*{You'll want to know: \sk{अदृष्टम्}}
\textbf{\sk{अदृष्टम्}} — \emph{adṛṣṭam} — \textquotedblleft{}the unseen; fate\textquotedblright{}.

\begin{quote}
\sk{यदि} \sk{अदृष्टं} \sk{साहाय्यं} \sk{करोति}, \sk{तर्हि} \sk{सः} \sk{जयति।}
\end{quote}

\begin{quote}
\emph{yadi adṛṣṭaṃ sāhāyyaṃ karoti, tarhi saḥ jayati.}
\end{quote}

\begin{quote}
If fate helps him, he wins.
\end{quote}

A neuter noun.

\begin{cousinweb}
\textbf{\sk{अदृष्टम्}} (\emph{adṛṣṭam}) is built on \emph{a}, \textquotedblleft{}not\textquotedblright{}, and \emph{dṛṣṭa}, \textquotedblleft{}seen\textquotedblright{}, from the root √dṛś, \textquotedblleft{}to see\textquotedblright{}: the unseen; MW gives \emph{adṛṣṭa} as what is beyond observation, the merit or demerit of one's deeds, fate (MW).
\end{cousinweb}

\begin{grammarlens}[title={What you've built}]
One more word for saying what happens if.
\end{grammarlens}

\subsection*{Your turn}
\begin{itemize}
\raggedright
  \item \emph{Say it:} \emph{adṛṣṭam}
  \item \emph{Say it:} \emph{adṛṣṭam}, once more
  \item \emph{Say it:} \emph{yadi adṛṣṭaṃ sāhāyyaṃ karoti, tarhi saḥ jayati.} — If fate helps him, he wins.
  \item \emph{Recall:} say the Sanskrit for a breaking, then the Sanskrit for one's word, then say \emph{adṛṣṭam} again
\end{itemize}

\begin{tcolorbox}[breakable,colback=teal!4,colframe=teal!35!black,title={Before you move on}]
What does \sk{अदृष्टम्} mean? (\textquotedblleft{}The unseen; fate\textquotedblright{}.) Say it once more.
\end{tcolorbox}
\section[\texorpdfstring{\sk{भवितव्यता} (bhavitavyatā)}{भवितव्यता (bhavitavyatā)}]{\sk{भवितव्यता} (bhavitavyatā) — what must be; destiny}

\label{lesson:SA-C536-bhavitavyata}



\begin{quote}
Before the new one: say the Sanskrit for one's word, then the Sanskrit for the unseen.
\end{quote}

\subsection*{You'll want to know: \sk{भवितव्यता}}
\textbf{\sk{भवितव्यता}} — \emph{bhavitavyatā} — \textquotedblleft{}what must be; destiny\textquotedblright{}.

\begin{quote}
\sk{यदि} \sk{भवितव्यता} \sk{एवम्} \sk{अस्ति}, \sk{तर्हि} \sk{किं} \sk{करोमि}?
\end{quote}

\begin{quote}
\emph{yadi bhavitavyatā evam asti, tarhi kiṃ karomi?}
\end{quote}

\begin{quote}
If that is what must be, what can I do?
\end{quote}

A feminine noun.

\begin{cousinweb}
\textbf{\sk{भवितव्यता}} (\emph{bhavitavyatā}) is built on \emph{bhavitavya}, \textquotedblleft{}what is to be\textquotedblright{}, from the root √bhū, \textquotedblleft{}to be\textquotedblright{}, with the abstract ending \emph{-tā}; MW gives \emph{bhavitavyatā} as inevitable consequence, necessity, destiny (MW).
\end{cousinweb}

\begin{grammarlens}[title={What you've built}]
One more word for saying what happens if.
\end{grammarlens}

\subsection*{Your turn}
\begin{itemize}
\raggedright
  \item \emph{Say it:} \emph{bhavitavyatā}
  \item \emph{Say it:} \emph{bhavitavyatā}, once more
  \item \emph{Say it:} \emph{yadi bhavitavyatā evam asti, tarhi kiṃ karomi?} — If that is what must be, what can I do?
  \item \emph{Recall:} say the Sanskrit for one's word, then the Sanskrit for the unseen, then say \emph{bhavitavyatā} again
\end{itemize}

\begin{tcolorbox}[breakable,colback=teal!4,colframe=teal!35!black,title={Before you move on}]
What does \sk{भवितव्यता} mean? (\textquotedblleft{}What must be; destiny\textquotedblright{}.) Say it once more.
\end{tcolorbox}
\section[\texorpdfstring{\sk{संयोगः} (saṃyogaḥ)}{संयोगः (saṃyogaḥ)}]{\sk{संयोगः} (saṃyogaḥ) — a coincidence; a conjunction}

\label{lesson:SA-C536-samyogah}



\begin{quote}
Before the new one: say the Sanskrit for the unseen, then the Sanskrit for what must be.
\end{quote}

\subsection*{You'll want to know: \sk{संयोगः}}
\textbf{\sk{संयोगः}} — \emph{saṃyogaḥ} — \textquotedblleft{}a coincidence; a conjunction\textquotedblright{}.

\begin{quote}
\sk{यदि} \sk{संयोगः} \sk{भवति}, \sk{तर्हि} \sk{त्वां} \sk{श्वः} \sk{द्रक्ष्यामि।}
\end{quote}

\begin{quote}
\emph{yadi saṃyogaḥ bhavati, tarhi tvāṃ śvaḥ drakṣyāmi.}
\end{quote}

\begin{quote}
If chance brings us together, I'll see you tomorrow.
\end{quote}

A masculine noun.

\begin{cousinweb}
\textbf{\sk{संयोगः}} (\emph{saṃyogaḥ}) is built on \emph{sam}, \textquotedblleft{}together\textquotedblright{}, and \emph{yoga}, \textquotedblleft{}joining\textquotedblright{}, from the root √yuj; MW gives \emph{saṃyoga} as conjunction, connection, union; \textquotedblleft{}a coincidence\textquotedblright{} is its modern use (MW).
\end{cousinweb}

\begin{grammarlens}[title={What you've built}]
One more word for saying what happens if.
\end{grammarlens}

\subsection*{Your turn}
\begin{itemize}
\raggedright
  \item \emph{Say it:} \emph{saṃyogaḥ}
  \item \emph{Say it:} \emph{saṃyogaḥ}, once more
  \item \emph{Say it:} \emph{yadi saṃyogaḥ bhavati, tarhi tvāṃ śvaḥ drakṣyāmi.} — If chance brings us together, I'll see you tomorrow.
  \item \emph{Recall:} say the Sanskrit for the unseen, then the Sanskrit for what must be, then say \emph{saṃyogaḥ} again
\end{itemize}

\begin{tcolorbox}[breakable,colback=teal!4,colframe=teal!35!black,title={Before you move on}]
What does \sk{संयोगः} mean? (\textquotedblleft{}A coincidence; a conjunction\textquotedblright{}.) Say it once more.
\end{tcolorbox}
\section[\texorpdfstring{\sk{यदृच्छया} (yadṛcchayā)}{यदृच्छया (yadṛcchayā)}]{\sk{यदृच्छया} (yadṛcchayā) — by chance}

\label{lesson:SA-C536-yadrcchaya}



\begin{quote}
Before the new one: say the Sanskrit for what must be, then the Sanskrit for a coincidence.
\end{quote}

\subsection*{You'll want to know: \sk{यदृच्छया}}
\textbf{\sk{यदृच्छया}} — \emph{yadṛcchayā} — \textquotedblleft{}by chance\textquotedblright{}.

\begin{quote}
\sk{यदि} \sk{यदृच्छया} \sk{मार्गे} \sk{मित्रं} \sk{पश्यामि}, \sk{तर्हि} \sk{मोदे।}
\end{quote}

\begin{quote}
\emph{yadi yadṛcchayā mārge mitraṃ paśyāmi, tarhi mode.}
\end{quote}

\begin{quote}
If I happen to meet a friend on the road, I am glad.
\end{quote}

\begin{cousinweb}
\textbf{\sk{यदृच्छया}} (\emph{yadṛcchayā}) is built from \emph{yad}, \textquotedblleft{}what, whatever\textquotedblright{}, and \emph{ṛcchā}, as MW gives it; MW records \emph{yadṛcchā} as self-will, accident, chance, and \emph{yadṛcchayā} as by accident (MW).
\end{cousinweb}

\begin{grammarlens}[title={What you've built}]
One more word for saying what happens if.
\end{grammarlens}

\subsection*{Your turn}
\begin{itemize}
\raggedright
  \item \emph{Say it:} \emph{yadṛcchayā}
  \item \emph{Say it:} \emph{yadṛcchayā}, once more
  \item \emph{Say it:} \emph{yadi yadṛcchayā mārge mitraṃ paśyāmi, tarhi mode.} — If I happen to meet a friend on the road, I am glad.
  \item \emph{Recall:} say the Sanskrit for what must be, then the Sanskrit for a coincidence, then say \emph{yadṛcchayā} again
\end{itemize}

\begin{tcolorbox}[breakable,colback=teal!4,colframe=teal!35!black,title={Before you move on}]
What does \sk{यदृच्छया} mean? (\textquotedblleft{}By chance\textquotedblright{}.) Say it once more.
\end{tcolorbox}
\section[\texorpdfstring{\sk{सम्पद्यते} (sampadyate)}{सम्पद्यते (sampadyate)}]{\sk{सम्पद्यते} (sampadyate) — turns out well}

\label{lesson:SA-C536-sampadyate}



\begin{quote}
Before the new one: say the Sanskrit for a coincidence, then the Sanskrit for by chance.
\end{quote}

\subsection*{You'll want to know: \sk{सम्पद्यते}}
\textbf{\sk{सम्पद्यते}} — \emph{sampadyate} — \textquotedblleft{}turns out well\textquotedblright{}.

\begin{quote}
\sk{यदि} \sk{सर्वं} \sk{सम्यक्} \sk{सम्पद्यते}, \sk{तर्हि} \sk{वयम्} \sk{उत्सवं} \sk{करिष्यामः।}
\end{quote}

\begin{quote}
\emph{yadi sarvaṃ samyak sampadyate, tarhi vayam utsavaṃ kariṣyāmaḥ.}
\end{quote}

\begin{quote}
If everything turns out well, we will celebrate.
\end{quote}

Past (it …): \sk{कार्यं} \sk{सम्पन्नवत्}, \emph{kāryaṃ sampannavat}.

\begin{cousinweb}
\textbf{\sk{सम्पद्यते}} (\emph{sampadyate}) is built on \emph{sam}, \textquotedblleft{}together, well\textquotedblright{}, and the root √pad, \textquotedblleft{}to fall, to go\textquotedblright{}; MW gives \emph{sampad} as to fall out well, to turn out well, to succeed; its past participle is \emph{sampanna} (MW).
\end{cousinweb}

\begin{grammarlens}[title={What you've built}]
One more word for saying what happens if.
\end{grammarlens}

\subsection*{Your turn}
\begin{itemize}
\raggedright
  \item \emph{Say it:} \emph{sampadyate}
  \item \emph{Say it:} \emph{sampadyate}, once more
  \item \emph{Say it:} \emph{yadi sarvaṃ samyak sampadyate, tarhi vayam utsavaṃ kariṣyāmaḥ.} — If everything turns out well, we will celebrate.
  \item \emph{Recall:} say the Sanskrit for a coincidence, then the Sanskrit for by chance, then say \emph{sampadyate} again
\end{itemize}

\begin{tcolorbox}[breakable,colback=teal!4,colframe=teal!35!black,title={Before you move on}]
What does \sk{सम्पद्यते} mean? (\textquotedblleft{}Turns out well\textquotedblright{}.) Say it once more.
\end{tcolorbox}
